
\medskip

%Une commune décide de remplacer le traditionnel feu d'artifice du 14 juillet par un
%spectacle de drones lumineux.
%
%\smallskip
%
%Pour le pilotage des drones, l'espace est muni d'un repère orthonormé \Oijk{} dont l'unité est la centaine de mètres.
%
%\smallskip
%
%La position de chaque drone est modélisée par un point et chaque drone est envoyé d'un point de départ D de coordonnées $(2~;~5~;~1)$.
%
%\smallskip
%
%On souhaite former avec des drones des figures en les positionnant dans un même plan $\mathcal{P}$.
%
%\smallskip
%
%Trois drones sont positionnés aux points A$(-1~;~-1~;~17)$, B$(4~;~-2~;~4)$ et C$(1~;~-3~;~7)$.

\begin{enumerate}
	\item %Justifier que les points A, B et C ne sont pas alignés.
On a $\vect{\text{AB}}\begin{pmatrix}5\\-1\\- 13\end{pmatrix}$ et $\vect{\text{AC}}\begin{pmatrix}2\\-2\\- 10\end{pmatrix}$.

Ces vecteurs ne sont pas colinéaires donc les points A, B et C  sont distincts et non alignés : ils définissent le plan $\mathcal{P} = $ (ABC).
\end{enumerate}

Dans la suite, on considère le vecteur $\vect{n} \begin{pmatrix}2\\-3\\1\end{pmatrix}$.

\begin{enumerate}[resume]
	\item 
	\begin{enumerate}
		\item %Justifier que $\vect{n}$ est normal au plan (ABC).
On a $\vect{n} \cdot \vect{\text{AB}} = 10  + 3 - 13 = 0$ et  $\vect{n} \cdot \vect{\text{AB}} = 4  + 6 - 10 = 0$.

Le vecteur $\vect{n}$ orthogonal à deux vecteurs non colinéaires du plan (ABC) est normal à ce plan
		\item %Démontrer qu'une équation cartésienne du plan $\mathcal{P}$ est $2x - 3y + z - 18 = 0$.
On sait qu'alors :
		
$M(x~;~y~;~z) \in  (\text{ABC}) \iff 2x - 3y + 1z + d = 0, \: d \in \R$.
		
Par exemple A$(-1~;~-1~;~17)\in (\text{ABC}) \iff - 2 + 3 + 17 + d = 0 \iff d = -18$, donc 

$M(x~;~y~;~z) \in  (\text{ABC}) \iff 2x - 3y + z - 18 = 0$.
\end{enumerate}
	\item %Le pilote des drones décide d'envoyer un quatrième drone en prenant comme trajectoire la droite $d$ dont une représentation paramétrique est donnée par 

\[ d~:~\left\{\begin{array}{l c l}
x&=&3t+2\\y&=&\phantom{3}t+5\\z&=&4t+1\end{array}\right.\text{, avec } t \in \R. \]
	\begin{enumerate}
		\item %Déterminer un vecteur directeur de la droite $d$.
On sait que les coordonnées d'un vecteur directeur $\vect{d}$ de $d$ sont les coefficients de $t$, soit $\vect{d}\begin{pmatrix}3\\1\\4\end{pmatrix}$
		\item %Afin que ce nouveau drone soit également placé dans le plan $\mathcal{P}$, déterminer par le calcul les coordonnées du point E, intersection de la droite $d$ avec le plan $\mathcal{P}$.
		Les coordonnées de E commun à $d$ et au plan (ABC) vérifient les équations paramétriques de $d$ et l'équation cartésienne de (ABC) soit le système
		
		$\left\{\begin{array}{l c l}
x&=&3t+2\\y&=&\phantom{3}t+5\\z&=&4t+1\\
2x - 3y + 1z - 18 &=& 0
\end{array}\right.\text{, avec } t \in \R.$

En remplaçant $x, \:y$ et $z$ par leurs valeurs en fonction de $t$ dans la dernière équation on obtient :

$2(3t + 2)- 3(t + 5) + 4t + 1 - 18 = 0 \iff 6t + 4 - 3y - 15 + 4t + 1 - 18 = 0\iff $

$7t - 28 = 0 \iff 7t = 28 \iff t = 4$.

D'où les coordonnées de E en remplaçant $t$ par dans les équations de $d$ : E(14~;~9~;~17)
	\end{enumerate}
	\item %Le pilote des drones décide d'envoyer un cinquième drone le long de la droite $\Delta$ qui passe par le point $D$ et qui est perpendiculaire au plan $\mathcal{P}$. Ce cinquième drone est placé lui aussi dans le plan $\mathcal{P}$, soit à l'intersection entre la droite $\Delta$ et le plan $\mathcal{P}$. On admet que le point F$(6~;~-1~;~3)$ correspond à cet emplacement.
La droite $\Delta$ contient D et a pour vecteur directeur un vecteur normal au plan (ABC) soit le vecteur $\vect{n}$.

Une équation de $\Delta$ est donc $\left\{\begin{array}{l c l}
x&=&2t+2\\y&=&-3t+5\\z&=&t+1\\
2x - 3y + 1z - 18 &=& 0
\end{array}\right.\text{, avec } t \in \R.$

En résolvant le système obtenu en rajoutant l'équation de $\mathcal{P}$ permet d'obtenir $t = 2$, d'où F$(6~;~-1~;~3)$.

Puisque la droite (DF) est perpendiculaire au plan $\mathcal{P}$, la distance DF est la (plus courte) distance du point au plan $\mathcal{P}$.

Or $\vect{\text{DF}}\begin{pmatrix}4\\-6\\2\end{pmatrix}$, d'où $\left\|\vect{\text{DF}}\right\|^2 = \text{DF}^2 = 4^2 + (-6)^2 + 2^2 = 16 + 36 + 4 = 56 = 4 \times 14$.

On a donc DF $ = \sqrt{4 \times 14} = \sqrt 4 \times \sqrt{14} = 2\sqrt{14}$ (en centaines de metres).
	
%Démontrer que la distance entre le point de départ D et le plan $\mathcal{P}$ vaut $2\sqrt{14}$ centaines de mètres.
	\item %L'organisatrice du spectacle demande au pilote d'envoyer un nouveau drone dans le plan (peu importe sa position dans le plan), toujours à partir du point D.
	
%Sachant qu'il reste $40$~secondes avant le début du spectacle et que le drone vole en trajectoire rectiligne à 18,6~$\text{m.s}^{-1}$, le nouveau drone peut-il arriver à temps ?
La plus petite distance du point de lancer au plan (ABC) est égale à $2\sqrt{14}$ centaines de mètres : il faut donc calculer le temps mis à la vitesse de 18,6~$\text{m.s}^{-1}$ par un drone pour effectuer ce parcours de D à F.

On sait que $v = \dfrac{d}{t}$, ou $t = \dfrac{d}{v} = \dfrac{\text{DF}}{v} = \dfrac{2\sqrt{14} \times 100}{18,6} \approx 40,23$ (secondes).

Conclusion : le drone n'arrivera pas à temps.
\end{enumerate}

